\let\oldprintchaptertitle=\printchaptertitle
\renewcommand{\printchaptertitle}[1]{%
	\vspace*{-75pt}
	\oldprintchaptertitle{#1}
}%
\myChapterStar{Résumé}{}{section}
\let\printchaptertitle=\oldprintchaptertitle

\noindent\textbf{Titre :} Ingénierie des applications web pilotées par des modèles BPMN

Ce mémoire s'intéresse à la génération d'applications web métier à partir de modèles de processus. Il part du constat que, dans de nombreuses organisations, les processus décrivent déjà une part essentielle du travail à informatiser : les activités à enchaîner, les acteurs concernés, les informations manipulées et les décisions à prendre. Pourtant, au moment de développer une application web, ces modèles restent souvent cantonnés à un rôle documentaire et leur contenu doit être réinterprété manuellement par les développeurs. Cette séparation fragilise la cohérence entre le besoin métier et le logiciel produit, complique la maintenance et rend l'évolution des applications plus coûteuse lorsque les procédures changent. L'objectif de ce travail est donc de concevoir un framework capable d'exploiter plus directement un modèle BPMN enrichi pour produire un socle applicatif web métier structuré, traçable et adaptable. Pour cela, nous avons conçu ProOWeb, un framework qui associe le modèle BPMN, chargé de représenter le déroulement du processus, à une spécification complémentaire. Cette spécification apporte les informations que le modèle ne porte pas suffisamment, notamment les rôles, les affectations, les données, les entrées, les sorties, les règles de visibilité et les indications nécessaires à la génération des formulaires. Le framework valide ces éléments, les organise dans un modèle pivot, puis les projette vers les artefacts du socle applicatif : vues de traitement, logique applicative, structures de données, règles d'accès, suivi des instances et traces de génération. Le travail met également l'accent sur la régénération contrôlée, afin que le projet cible puisse évoluer sans rompre le lien entre le processus, la spécification et le code. Ainsi, ProOWeb ne vise pas seulement à automatiser l'écriture de fichiers, mais à rapprocher durablement la modélisation métier et la réalisation logicielle.

\vspace{0.8cm}

\noindent\textbf{Mots clés :} Processus métier, application web métier, BPMN, ProOWeb, génération d'application, ingénierie dirigée par les modèles, spécification enrichie, régénération contrôlée.

\myCleanStarChapterEnd
